\begin{pastexam}[label={px:gw3-111-ds-02}]{關務三等資訊處理 111 年 第二題}
（申論；20 分）

遞迴函數（recursive function）起始數值與遞迴關係定義為：\\
P(0) = P(1) = P(2) = 1,P(n) = P(n$-$1)-2P(n$-$2)+P(n$-$3), $\forall$n >=3
\begin{itemize}[leftmargin=2.2em]
\item[㈠] 請問 P(n)的前 5 個值依序為：1,1,1,及那兩個數字？（10 分）
\item[㈡] 請問 P(6)+P(8)的值為何？（10 分）
\end{itemize}
\end{pastexam}

\begin{pastsoln}
答案：㈠ $0$ 和 $-1$；㈡ $P(6) + P(8) = 2 + (-3) = -1$。

照定義由小到大逐項代入（資結書第~2.5~節的線性遞迴；這題只要求數值，直接算最穩）：
\begin{center}
\small
\begin{tabular}{cll}
\toprule
$n$ & $P(n) = P(n-1) - 2P(n-2) + P(n-3)$ & $P(n)$ \\
\midrule
0, 1, 2 & 起始值 & 1, 1, 1 \\
3 & $P(2) - 2P(1) + P(0) = 1 - 2 + 1$ & $0$ \\
4 & $P(3) - 2P(2) + P(1) = 0 - 2 + 1$ & $-1$ \\
5 & $P(4) - 2P(3) + P(2) = -1 - 0 + 1$ & $0$ \\
6 & $P(5) - 2P(4) + P(3) = 0 + 2 + 0$ & $2$ \\
7 & $P(6) - 2P(5) + P(4) = 2 - 0 - 1$ & $1$ \\
8 & $P(7) - 2P(6) + P(5) = 1 - 4 + 0$ & $-3$ \\
\bottomrule
\end{tabular}
\end{center}

\textbf{㈠} 前 $5$ 個值是 $P(0..4) = 1, 1, 1, \mathbf{0}, \mathbf{-1}$，所求的兩個數字是 $0$ 和 $-1$。

\textbf{㈡} $P(6) + P(8) = 2 + (-3) = \mathbf{-1}$。

注意遞迴式的中間項是 $-2P(n-2)$，前一項是負數時要變號（例如 $P(6)$ 裡的 $-2P(4) = -2 \cdot (-1) = +2$），這是最常
算錯的地方。 若寫成程式，直接照定義遞迴會重複計算（呼叫次數隨 $n$ 指數成長），應該像上表一樣由小到大用陣列或
三個變數往前推，$O(n)$ 時間。
\end{pastsoln}
